\section{The Restless Heart and the Kingdom Prepared}

Every autumn, before the Church clothes her altars in black for All Souls' Day or opens the last Sunday of the year upon the darkened sun and falling stars, she first lifts the veil upon the end for which all things were made. On the first of November, vested in white, she enters the sanctuary singing the Introit of All Saints: \textit{Gaudeamus omnes in Domino, diem festum celebrantes sub honore Sanctorum omnium: de quorum solemnitate gaudent Angeli, et collaudant Filium Dei}---``Let us all rejoice in the Lord, celebrating a festival day under the honor of all the Saints: on whose solemnity the Angels rejoice, and praise the Son of God.''\endnote{1962 \textit{Missale Romanum}, In Festo Omnium Sanctorum (1 November), Introit, collated with the 1922 Tours Mame typical edition.} The last word of Catholic eschatology is not dread, nor mere survival after death, but the King's summons to those at His right hand: ``Come, ye blessed of my Father, possess you the kingdom prepared for you from the foundation of the world'' (Matt 25:34).

Christ Himself sharpens the contrast between the two verdicts of Matthew 25. The everlasting fire was ``prepared for the devil and his angels'' (Matt 25:41), an exile men enter only by sharing the fall of the rebel spirits. The kingdom was ``prepared for you from the foundation of the world'' (Matt 25:34). Man was fashioned for heaven. The thirst for truth that no science exhausts, the ache for beauty that fades when grasped, and the hunger for love untouched by death all bear the mark of that destination.

St.\ Augustine gives that longing its enduring voice in the opening lines of his \textit{Confessions}:

\begin{quote}
Great art Thou, O Lord, and greatly to be praised; great is Thy power, and of Thy wisdom there is no end. And man, being a part of Thy creation, desires to praise Thee, man, who bears about with him his mortality, the witness of his sin, even the witness that Thou ``resistest the proud,''---yet man, this part of Thy creation, desires to praise Thee. Thou movest us to delight in praising Thee; for Thou hast formed us for Thyself, and our hearts are restless till they find rest in Thee.\endnote{St.\ Augustine, \textit{Confessions} I.1.1, trans.\ J.\ G.\ Pilkington, NPNF1 vol.\ 1.}
\end{quote}

Having tested rhetorical renown, ambition, pleasure, and imperial favor, Augustine saw that restlessness is not a flaw in human nature; it is the tremor of a spirit shaped for an infinite Object until it locks upon its pole.

Long before the Incarnation, the Psalter sang out of that same desire. ``As the hart panteth after the fountains of water; so my soul panteth after thee, O God. My soul hath thirsted after the strong living God; when shall I come and appear before the face of God?'' (Ps 41:2--3).\endnote{Scripture quotations and verse numbering throughout follow the Challoner revision of the Douay-Rheims Bible and the Clementine Vulgate.} The psalmist asks not for tranquil gardens, but for a Face: ``For what have I in heaven? and besides thee what do I desire upon earth? For thee my flesh and my heart hath fainted away: thou art the God of my heart, and the God that is my portion for ever'' (Ps 72:25--26). Again he cries out for the King's courts: ``How lovely are thy tabernacles, O Lord of hosts! my soul longeth and fainteth for the courts of the Lord. \dots\ Blessed are they that dwell in thy house, O Lord: they shall praise thee for ever and ever. \dots\ For better is one day in thy courts above thousands'' (Ps 83:2--3, 5, 11). And Psalm 16 gathers all longing into one line: ``But as for me, I will appear before thy sight in justice: I shall be satisfied when thy glory shall appear'' (Ps 16:15).

In the Gospels, Jesus Christ reveals Himself as the way to the Father's house and the door of entry. At the Last Supper He steadies His Apostles with plain assurance: ``In my Father's house there are many mansions. If not, I would have told you: because I go to prepare a place for you. And if I shall go and prepare a place for you, I will come again and will take you to myself: that where I am, you also may be'' (John 14:2--3). In His High Priestly Prayer He names the substance of that life: ``Now this is eternal life: That they may know thee, the only true God, and Jesus Christ, whom thou hast sent'' (John 17:3); and He utters His sovereign testament: ``Father, I will that where I am, they also whom thou hast given me may be with me: that they may see my glory which thou hast given me, because thou hast loved me before the creation of the world'' (John 17:24). In the parable of the talents, the reward is admission into the Master's own bliss: ``Enter thou into the joy of thy lord'' (Matt 25:21); and from the Cross He promises the penitent thief immediate entry: ``Amen I say to thee: This day thou shalt be with me in paradise'' (Luke 23:43).

The Apostles measure every earthly hope against that height. St.\ Paul declares that ``eye hath not seen, nor ear heard: neither hath it entered into the heart of man, what things God hath prepared for them that love him'' (1 Cor 2:9), yet contrasts our present twilight with unveiled sight: ``We see now through a glass in a dark manner: but then face to face. Now I know in part: but then I shall know even as I am known'' (1 Cor 13:12). St.\ John joins vision directly to likeness: ``Dearly beloved, we are now the sons of God: and it hath not yet appeared what we shall be. We know that when he shall appear we shall be like to him: because we shall see him as he is'' (1 John 3:2).

Why can nothing less than seeing God appease the human heart? St.\ Thomas Aquinas devotes the opening five questions of the \textit{Prima Secundae} to answering that question with quiet rigor. Across the first seven articles of Question 2 he tests every created rival in turn. Wealth fails because money is sought only for other things and sharpens covetousness (\textit{ST} I-II, q.\ 2, a.\ 1); honor fails because it lies in fickle givers (a.\ 2); fame fails because human opinion errs (a.\ 3); power fails because it can serve evil as readily as good (a.\ 4); bodily health and beauty fail because the body exists for the soul (a.\ 5); sensible pleasure fails because delight follows a good rather than constituting it (a.\ 6); even created goods of the soul---knowledge and virtue---cannot be the last end, because the soul's powers reach beyond themselves (a.\ 7).

In Article 8, Aquinas gathers every refutation into one principle:

\begin{quote}
It is impossible for any created good to constitute man's happiness. For happiness is the perfect good, which lulls the appetite altogether; else it would not be the last end, if something yet remained to be desired. Now the object of the will, i.e.\ of man's appetite, is the universal good; just as the object of the intellect is the universal true. Hence it is evident that naught can lull man's will, save the universal good. This is to be found, not in any creature, but in God alone; because every creature has goodness by participation. Wherefore God alone can satisfy the will of man, according to the words of Ps.\ 102:5: ``Who satisfieth thy desire with good things.'' Therefore God alone constitutes man's happiness.\endnote{St.\ Thomas Aquinas, \textit{Summa Theologiae} I-II, q.\ 2, a.\ 8, trans.\ Fathers of the English Dominican Province.}
\end{quote}

Even the entire created cosmos, joined with every angel, remains a finite, participated good (\textit{ST} I-II, q.\ 2, a.\ 8, ad 2--3). Nor can philosophical knowledge of God through created effects suffice (q.\ 3, a.\ 8). Just as seeing an eclipse drives the mind to seek the essence of its cause, ``if therefore the human intellect, knowing the essence of some created effect, knows no more of God than `that He is'; the perfection of that intellect does not yet reach simply the First Cause, but there remains in it the natural desire to seek the cause. Wherefore it is not yet perfectly happy. Consequently, for perfect happiness the intellect needs to reach the very Essence of the First Cause.''\endnote{St.\ Thomas Aquinas, \textit{Summa Theologiae} I-II, q.\ 3, a.\ 8.}

Once attained, that happiness can never be lost (\textit{ST} I-II, q.\ 5, a.\ 4). Against Origen's conjecture that blessed souls might grow sated and fall away, Aquinas observes that fear of loss---or false security---would poison beatitude itself. More fundamentally, the will turns away from a good only when it sees some defect in it or some rival outside it; before the infinite Essence of God, where every good is present without shadow, the will cleaves to Him with joyful, unshakeable necessity.

How, then, does the Christian spend the years of exile before that vision dawns? Preaching on 1 John 3:2, St.\ Augustine turns metaphysics directly into spiritual discipline:

\begin{quote}
The whole life of a good Christian is an holy desire. Now what thou longest for, thou dost not yet see: howbeit by longing, thou art made capable, so that when that is come which thou mayest see, thou shall be filled. For just as, if thou wouldest fill a bag, and knowest how great the thing is that shall be given, thou stretchest the opening of the sack or the skin, or whatever else it be; thou knowest how much thou wouldest put in, and seest that the bag is narrow; by stretching thou makest it capable of holding more: so God, by deferring our hope, stretches our desire; by the desiring, stretches the mind; by stretching, makes it more capacious.\endnote{St.\ Augustine, \textit{Homilies on the First Epistle of John} IV.6, trans.\ H.\ Browne, NPNF1 vol.\ 7.}
\end{quote}

Delay is enlargement; yet stretching the vessel requires emptying it as well. ``Suppose that God would fill thee with honey,'' Augustine warns his flock in the next sentence: ``if thou art full of vinegar, where wilt thou put the honey? That which the vessel bore in it must be poured out: the vessel itself must be cleansed; must be cleansed, albeit with labor, albeit with hard rubbing, that it may become fit for that thing, whatever it be.''
